\documentclass[12pt,a4paper]{article}
\usepackage[utf8]{inputenc}
\usepackage[T1]{fontenc}
\usepackage[brazil]{babel}
\usepackage{indentfirst}
\usepackage[margin=2.5cm]{geometry}

\title{Resenha Cr\'itica: Artigo 5 ---\\
Software Architecture: A Roadmap}
\author{Gustavo Pessoa Firmino Duarte}
\date{29 de Mar\c{c}o de 2026}

\begin{document}
\maketitle

\section{Introdu\c{c}\~ao}

Publicado em 2000 nos anais do \textit{International Conference on Software
Engineering} (ICSE), o artigo ``Software Architecture: A Roadmap'', de David
Garlan, tra\c{c}a um panorama da \'area de arquitetura de software ao fim do
s\'eculo XX e aponta os desafios que deveriam guiar a pesquisa no s\'eculo
seguinte. Escrito no contexto de uma s\'erie sobre o futuro da engenharia de
software, o trabalho parte de conquistas consolidadas --- estilos arquiteturais,
linguagens de descri\c{c}\~ao de arquitetura (ADLs), an\'alise de trade-offs ---
para identificar lacunas ainda abertas: a falta de ferramentas de suporte
maduras, a dificuldade de conectar arquitetura ao c\'odigo e a necessidade de
m\'etodos de an\'alise quantitativa de qualidades. Vinte e cinco anos depois,
boa parte dessas lacunas permanece relevante.

\section{Estilos Arquiteturais e Desafios do Campo}

Garlan define arquitetura de software como o conjunto de componentes, conectores
e a configura\c{c}\~ao que os combina, regida por propriedades observ\'aveis ao
n\'ivel do sistema como um todo. Ele descreve o avanço na cataloga\c{c}\~ao de
\textit{estilos arquiteturais} --- pipe-and-filter, cliente-servidor, em camadas
(\textit{layered}), orientado a eventos, reposit\'orio e espa\c{c}o de dados
compartilhado, entre outros. Cada estilo captura uma solu\c{c}\~ao recorrente
para um conjunto de problemas e carrega impl\'icitas suas vantagens e limitações.

O artigo aponta que, embora estilos sejam bem compreendidos teoricamente, a
pr\'atica de doc\'umenta\c{c}\~ao arquitetural permanecia informal e inconsistente.
As ADLs existentes (Aesop, Rapide, Wright) eram acad\^emicas e sem ado\c{c}\~ao
industrial significativa. Garlan tamb\'em destaca a tens\~ao entre vis\~oes
m\'ultiplas de uma arquitetura --- a vis\~ao l\'ogica, a vis\~ao de processo e a
vis\~ao f\'isica (antecipando o modelo 4+1 de Kruchten) --- e a dificuldade de
manter a consist\^encia entre essas vis\~oes conforme o sistema evolui.

Entre os desafios que o autor elenca para o futuro, destacam-se: (1) ferramentas
que tornem a arquitetura document\'avel e verific\'avel automaticamente;
(2) m\'etodos para analisar qualidades de sistema (desempenho, disponibilidade,
segurança) diretamente na representa\c{c}\~ao arquitetural; (3) suporte ao ciclo
de vida completo, conectando decis\~oes arquiteturais \`a implementa\c{c}\~ao e
\`a opera\c{c}\~ao. Esses desafios moldaram pesquisas nas d\'ecadas seguintes,
incluindo o modelo C4, o ATAM (\textit{Architecture Tradeoff Analysis Method})
e os microsservi\c{c}os.

\section{Conex\~ao com a Pr\'atica Profissional}

A leitura deste artigo me fez refletir sobre como a arquitetura de software \'e
tratada nos projetos em que participei. No meu est\'agio na ArcelorMittal, os
sistemas de suporte internacional possu\'iam arquiteturas impl\'icitas ---
camadas e componentes que existiam na mente das pessoas mais experientes da equipe,
mas raramente estavam documentadas de forma expl\'icita. A onboarding de novos
membros dependia de convers\'ation, n\~ao de diagramas. Esse cenário ilustra
precisamente o problema que Garlan identifica: sem representa\c{c}\~oes arquiteturais
formais e acess\'iveis, o conhecimento fica concentrado em pessoas espec\'ificas
e o risco de erros arquiteturais n\~ao detectados aumenta.

No meu MVP de e-commerce, tomei decis\~oes arquiteturais intuitivas --- separar
rotas, servi\c{c}os e modelos --- mas nunca as documentei explicitamente. Ap\'os
ler Garlan, percebo que tomei decis\~oes de estilo arquitetural (uma arquitetura
em camadas com API RESTful) sem articular os trade-offs envolvidos: por que n\~ao
um BFF (\textit{Backend for Frontend})? Por que n\~ao event-driven? Essas escolhas
ficaram impl\'icitas no c\'odigo, exatamente o antipadr\~ao que o autor critica.

\section{Conclus\~ao}

O ``Roadmap'' de Garlan permanece surpreendentemente atual. Os desafios que ele
elencou em 2000 --- ferramentas de documenta\c{c}\~ao, an\'alise quantitativa de
qualidades, conformidade arquitetural --- foram parcialmente endereçados por
iniciativas como o modelo C4 (Simon Brown), o TOGAF e pr\'aticas de
\textit{architecture fitness functions} (Ford et al.). Mas a br\'echa entre
arquitetura e implementa\c{c}\~ao, a ``d\'ervia arquitetural'', continua sendo
um dos problemas mais comuns na ind\'ustria. O artigo serve como um lembrete de
que arquitetura n\~ao \'e apenas um diagramas inicial --- \'e uma pr\'atica
cont\'inua que precisa de ferramentas, m\'etodos e cultura organizacional para
ser efetiva.

\end{document}
